\documentclass[12pt]{article}
\usepackage[utf8]{inputenc}
\usepackage[T1]{fontenc}
\usepackage{hyperref}
\title{Sunwar Language Resources for K-12 Curriculum Development}
\author{Compiled by Claude Sonnet 5.5 (claude-sonnet-5-5) for Andr\'as Kornai}
\date{}
\begin{document}
\maketitle

\section{Introduction}
Sunwar (Sunuwar; ISO 639-3 suz; Glottocode sunw1242) is a Sino-Tibetan (Kiranti) language of eastern Nepal, spoken in Okhaldhunga and Ramechhap districts. Speaker counts differ between sources. The task description gives about 38,000, and the Endangered Languages Project entry gives about 25,000. The searches found a grammar and dictionaries but no school materials.

\section{Primary Educational Resources}
No primer, textbook or school curriculum was found in the search results, and nothing on mother-tongue schooling.

\section{Linguistic and Reference Materials}
\begin{itemize}
\item Borchers, \emph{A Grammar of Sunwar: Descriptive Grammar, Paradigms, Texts and Glossary} (Brill, 2008): \url{https://www.universiteitleiden.nl/en/research/research-output/humanities/a-grammar-of-sunwar.-descriptive-grammar-paradigms-texts-and-glossary}
\item Leiden repository copy of the same work: \url{https://scholarlypublications.universiteitleiden.nl/access/item%3A2920095/view}
\item Sunuwar dictionary on Webonary (Bieri and Schulze, with Sunuwar collaborators): \url{https://www.webonary.work/sunuwar/overview/preface/}
\item Webonary credits and acknowledgements page: \url{https://www.webonary.work/sunuwar/?p=186002}
\item Bieri and Schulze, \emph{A Vocabulary of the Sunwar Language} (1971), bibliographic record: \url{https://glottolog.org/resource/reference/id/154287}
\end{itemize}

\section{Religious and Community Resources}
None found in the search results.

\section{Digital and Audio Resources}
The Webonary dictionary site is the main digital resource found. No audio or video was found. Endangered Languages Project entry: \url{https://endangeredlanguages.com/node/97180}

\section{Curriculum Development Context}
The grammar and online dictionary give a base for materials. Orthography and literacy materials, and any community organisations, should be confirmed with Sunuwar speakers, since none appeared in these searches.
\end{document}